\subsection{COMPLETION OF A TOPOLOGICAL RING}

When we speak of the uniformity of a topological ring, we always mean the uniformity of its additive group, unless the contrary is expressly stated; in particular, A is said to be a complete ring if the additive group of A is complete.

Let A be a Hausdorff topological ring; as an additive group A can be considered as a dense subgroup of a complete Hausdorff commutative group \(\hat{\mathbf{A}}\) , which is determined up to isomorphism (§ 3, no. 5, Theorem 2). In order that we should be able to consider A as a subring of a complete ring, it is necessary to be able to extend the function \(xy\) by continuity to the space \(\hat{\mathbf{A}} \times \hat{\mathbf{A}}\) . The possibility of this extension will follow from the following more general theorem:

THEOREM 1. Let E, F, G be three complete Hausdorff commutative groups; let A be a dense subgroup of E and let B be a dense subgroup of F. If f is a Z continuous Z-bilinear (*) mapping of A × B into G, then f can be extended by continuity to a continuous Z-bilinear mapping of E × F into G.

Let \((x_{0}, y_{0})\) be an arbitrary point of \(E \times F\) , and let \(\mathfrak{U}\) , \(\mathfrak{B}\) be the traces on A, B respectively of the neighbourhood filters of \(x_{0}, y_{0}\) ( \(\mathfrak{U}\) , \(\mathfrak{B}\) are filters, by hypothesis). To show that f can be extended by continuity, it is enough to show that \(f(\mathfrak{U} \times \mathfrak{V})\) is a Cauchy filter-base on G (Chapter II, § 3, no. 6, Proposition 11). Consider the identity

\[
f \left(x ^ {\prime}, y ^ {\prime}\right) - f (x, y) = f \left(x ^ {\prime} - x, y _ {1}\right) + f \left(x _ {1}, y ^ {\prime} - y\right) + f \left(x ^ {\prime} - x, y ^ {\prime} - y _ {1}\right).
\]

We shall show that by taking \((x, y)\) and \((x', y')\) in a sufficiently small set of \(\mathfrak{U} \times \mathfrak{B}\) , and by choosing \(x_{1}\) and \(y_{1}\) suitably, we can make each of the terms on the right-hand side very small. Let W be any neighbourhood of o in G; since f is continuous at \((o, o) \in A \times B\) , there is a set \(U \in \mathfrak{U}\) and a set \(V \in \mathfrak{B}\) such that \(f(x' - x, y' - y) \in W\) whenever \(x \in U\) , \(x' \in U\) , \(y \in V\) , \(y' \in V\) . Take a point \(x_{1} \in U\) and a point \(y_{1} \in V\) ; then for all \(x, x'\) in U and all \(y, y'\) in V we shall have

\[
f \left(x ^ {\prime} - x, y ^ {\prime} - y _ {1}\right) + f \left(x - x _ {1}, y ^ {\prime} - y\right) \in \mathrm{W} + \mathrm{W}.
\]

On the other hand, the partial mapping \(x \to f(x, y_1)\) is continuous on A; hence there is a set \(U' \subset U\) , belonging to \(\mathfrak{U}\) , and such that, whenever \(x\) and \(x'\) belong to \(U'\) , we have \(f(x' - x, y_1) \in W\) . Likewise there exists \(V' \subset V\) belonging to \(\mathfrak{B}\) such that, whenever \(y\) and \(y'\) belong to \(V'\) , we have \(f(x_1, y' - y) \in W\) . Consequently, if \((x, y)\) and \((x', y')\) are any two points of \(U' \times V'\) , then

\[
f \left(x ^ {\prime}, y ^ {\prime}\right) - f (x, y) \in W + W + W + W;
\]

this proves the existence of the extension \(\overline{f}\) of \(f\) . The fact that \(\overline{f}\) is \(\mathbf{Z}\) -bilinear is an immediate consequence of the principle of extension of identities (Chapter I, § 8, no. 1, Corollary 1 to Proposition 2).

Q.E.D.

In the application of this theorem to a Hausdorff topological ring A, we take E, F and G to be \(\hat{A}\) , A and B to be the ring A, and f to be the Z-bilinear mapping \((x, y) \to xy\) , which by hypothesis is continuous. We denote again by xy the value of the extended function on \(\hat{A} \times \hat{A}\) ; this function is a law of composition on \(\hat{A}\) , and to say that it is Z-bilinear means that it is distributive on both sides with respect to addition; and it is also associative, by the principle of extension of identities. Consequently:

PROPOSITION 6. A Hausdorff topological ring A is isomorphic to a dense subring of a complete Hausdorff ring \(\hat{\mathbf{A}}\) , which is determined up to isomorphism (and is called the completion of A).

If A is commutative (resp. has an identity element) then the same is true of \(\hat{A}\) (principle of extension of identities).

Let A be a topological ring, not necessarily Hausdorff; let N be the closure of \(\{o\}\) in A, and let \(A' = A/N\) be the Hausdorff ring associated with A. Then \(\hat{A}'\) , the completion of \(A'\) , is called the Hausdorff completion of A and is also denoted by \(\hat{A}\) . One shows as in §3, no. 4, Proposition 8 that every continuous ring homomorphism u of A with a complete Hausdorff topological ring C can be uniquely factorized as \(u = v \circ \varphi\) , where v is a continuous homomorphism of \(\hat{A}\) into C and \(\varphi\) is the canonical mapping of A into \(\hat{A}\) . If A, B are two topological rings, and \(u : A \to B\) is a continuous homomorphism, there is therefore a unique continuous homomorphism \(\hat{u} : \hat{A} \to \hat{B}\) such that the diagram

\[
\begin{array}{c c c} \mathbf {A} & \stackrel {{u}} {{\longrightarrow}} & \mathbf {B} \\ \varphi \Big \downarrow & & \Big \downarrow \psi \\ \hat {\mathbf {A}} & \stackrel {{\hat {u}}} {{\longrightarrow}} & \hat {\mathbf {B}} \end{array}
\]

is commutative ( \(\varphi\) and \(\psi\) being the canonical mappings); we have only to apply the preceding result to \(\psi \circ u\) .

